\documentclass[11pt, a4paper]{article}

% bfw-style.tex — gemeinsame Style-Definitionen für alle Handouts/Aufgabenhefte
% Voraussetzung: Kompilation mit xelatex (wegen fontspec)

\usepackage[ngerman]{babel}
\usepackage{geometry}
\geometry{a4paper, top=2.2cm, bottom=2.2cm, left=2.3cm, right=2.3cm}

\usepackage{fontspec}
% Fallback-Kette: Source Sans Pro -> Calibri -> Segoe UI -> Latin Modern Sans
% (Latin Modern ist immer mit MiKTeX vorhanden)
\IfFontExistsTF{Source Sans Pro}{%
  \setmainfont{Source Sans Pro}[Ligatures=TeX]%
}{\IfFontExistsTF{Calibri}{%
  \setmainfont{Calibri}[Ligatures=TeX]%
}{\IfFontExistsTF{Segoe UI}{%
  \setmainfont{Segoe UI}[Ligatures=TeX]%
}{%
  \setmainfont{Latin Modern Sans}[Ligatures=TeX]%
}}}
\IfFontExistsTF{Source Code Pro}{%
  \setmonofont{Source Code Pro}[Scale=0.92]%
}{\IfFontExistsTF{Consolas}{%
  \setmonofont{Consolas}[Scale=0.92]%
}{%
  \setmonofont{Latin Modern Mono}[Scale=0.92]%
}}

\usepackage{xcolor}
% Farbpalette: hell, freundlich, modern
\definecolor{bfwPrimary}{HTML}{0EA5A4}    % Petrol/Türkis - Primär
\definecolor{bfwPrimaryDark}{HTML}{0F766E} % dunkler Petrol für Headings
\definecolor{bfwAccent}{HTML}{F59E0B}     % Amber - Akzent
\definecolor{bfwSuccess}{HTML}{10B981}    % Grün - Erfolg / Lösung
\definecolor{bfwWarn}{HTML}{F97316}       % Orange - Achtung
\definecolor{bfwInk}{HTML}{0F172A}        % fast Schwarz
\definecolor{bfwInkSoft}{HTML}{475569}    % gedämpftes Grau
\definecolor{bfwBgSoft}{HTML}{F8FAFC}     % off-white
\definecolor{bfwBgTeal}{HTML}{ECFEFF}     % helles Türkis
\definecolor{bfwBgAmber}{HTML}{FEF3C7}    % helles Gelb
\definecolor{bfwBgRose}{HTML}{FFE4E6}     % helles Rosé für Eselsbrücken

\usepackage[colorlinks=true, linkcolor=bfwPrimaryDark, urlcolor=bfwPrimaryDark]{hyperref}
\usepackage{enumitem}
\setlist[itemize]{leftmargin=*, itemsep=2pt, topsep=2pt}
\setlist[enumerate]{leftmargin=*, itemsep=2pt, topsep=2pt}

\usepackage[most]{tcolorbox}
\usepackage{booktabs}
\usepackage{tikz}
\usetikzlibrary{decorations.pathmorphing, positioning, shapes.geometric, arrows.meta}

% Merkkästchen — wichtige Definition / Kernpunkt
\newtcolorbox{merksatz}[1][]{%
  enhanced, breakable,
  colback=bfwBgTeal,
  colframe=bfwPrimary,
  boxrule=0.6pt,
  arc=4pt,
  left=10pt, right=10pt, top=8pt, bottom=8pt,
  fonttitle=\bfseries\color{bfwPrimaryDark},
  title={\faLightbulb\ Merksatz},
  #1
}

% Eselsbrücke — Mnemoniks
\newtcolorbox{eselsbruecke}[1][]{%
  enhanced, breakable,
  colback=bfwBgRose,
  colframe=bfwAccent,
  boxrule=0.6pt,
  arc=4pt,
  left=10pt, right=10pt, top=8pt, bottom=8pt,
  fonttitle=\bfseries\color{bfwWarn},
  title={Eselsbrücke},
  #1
}

% Beispiel-Box
\newtcolorbox{beispiel}[1][]{%
  enhanced, breakable,
  colback=bfwBgSoft,
  colframe=bfwInkSoft,
  boxrule=0.4pt,
  arc=3pt,
  left=10pt, right=10pt, top=8pt, bottom=8pt,
  fonttitle=\bfseries\color{bfwInk},
  title={Beispiel},
  #1
}

% Lösungs-Box (für Aufgabenhefte)
\newtcolorbox{loesung}[1][]{%
  enhanced, breakable,
  colback=bfwBgAmber!40,
  colframe=bfwSuccess,
  boxrule=0.6pt,
  arc=3pt,
  left=10pt, right=10pt, top=8pt, bottom=8pt,
  fonttitle=\bfseries\color{bfwSuccess!70!black},
  title={Lösung},
  #1
}

% Glossar-Eintrag
\newcommand{\glossareintrag}[2]{%
  \par\noindent\textbf{\color{bfwPrimaryDark}#1} \enspace #2\par\smallskip
}

% Section-Stil — etwas farbiger als Standard
\usepackage{titlesec}
\titleformat{\section}
  {\Large\bfseries\color{bfwPrimaryDark}}
  {\thesection}{0.6em}{}
\titleformat{\subsection}
  {\large\bfseries\color{bfwPrimary}}
  {\thesubsection}{0.5em}{}
\titleformat{\subsubsection}
  {\normalsize\bfseries\color{bfwInk}}
  {\thesubsubsection}{0.4em}{}

% TikZ-Stil "skizzenhaft" — etwas weicher als Rough.js, aber stabil
\tikzset{
  roughbox/.style = {
    draw=bfwPrimaryDark, thick, fill=bfwBgTeal,
    rounded corners=4pt, inner sep=6pt
  },
  roughaccent/.style = {
    draw=bfwAccent, thick, fill=bfwBgAmber,
    rounded corners=4pt, inner sep=6pt
  },
  roughline/.style = {
    draw=bfwInkSoft, thick, -{Stealth[length=2.5mm]}
  }
}

% Bullet-Symbole (bewusst kein FontAwesome — vermeidet Font-Installations-Probleme)
\newcommand{\faLightbulb}{$\bullet$}
\newcommand{\faBookmark}{$\bullet$}

% Header/Footer
\usepackage{fancyhdr}
\pagestyle{fancy}
\fancyhf{}
\renewcommand{\headrulewidth}{0pt}
\fancyfoot[L]{\small\color{bfwInkSoft}\thepage}
\fancyfoot[R]{\small\color{bfwInkSoft} BFW M-064 Beschaffung im Büromanagement}


\title{\vspace*{-1.5cm}\Huge\bfseries\color{bfwPrimaryDark}Aufgabenheft Tag~5\\[0.4em]
  \large\color{bfwInkSoft}\textnormal{Annahme und Lagerung der Waren --- Übungen im IHK-Klausurformat}}
\author{\textsf{Beschaffung im Büromanagement}\\
\textsf{\small\copyright\ Can Siebert\,\textperiodcentered\,2026}}
\date{Selbstlernmaterial \textperiodcentered{} 18.05.2026}

\begin{document}
\maketitle
\thispagestyle{empty}

\begin{merksatz}[title={\bfseries So nutzen Sie dieses Aufgabenheft}]
Bearbeiten Sie alle Aufgaben zunächst \emph{ohne} in die Lösungen zu schauen.
Beachten Sie die \emph{Operatoren} und schreiben Sie Antworten in vollständigen Sätzen.
Lösungen ab Seite~\pageref{loesungen}.
\end{merksatz}

\tableofcontents
\vspace{1em}
\hrule

\section{Aufgaben}

\subsection*{Aufgabe~1 --- Wareneingangskontrolle und Mängelrüge (12 Punkte)}

\noindent\textbf{\color{bfwAccent}YouTube-Vorbereitung:}\enspace \href{https://www.youtube.com/watch?v=aMiVDStRLEU}{Die Wareneingangskontrolle (Prozubi)} (\url{https://www.youtube.com/watch?v=aMiVDStRLEU})

\medskip
Sie nehmen am 20.~Mai eine Lieferung von 100~Aktenordnern entgegen. Beim Auspacken
stellen Sie fest, dass nur 95~Stück geliefert wurden, und 8~davon sind verschmutzt
und teilweise eingerissen.

\begin{enumerate}[label=(\alph*)]
  \item \textbf{Beschreiben} Sie den Ablauf der Wareneingangskontrolle in 4~Schritten. \hfill (4~P.)
  \item \textbf{Unterscheiden} Sie offene und verdeckte Mängel. \hfill (2~P.)
  \item \textbf{Erläutern} Sie §\,377 HGB: Welche Pflicht hat der Käufer beim Handelskauf? \hfill (3~P.)
  \item \textbf{Verfassen} Sie kurz eine Mängelrüge an die \emph{Bürofix AG}. \hfill (3~P.)
\end{enumerate}

\subsection*{Aufgabe~2 --- Aufgaben der Lagerhaltung (8 Punkte)}

\begin{enumerate}[label=(\alph*)]
  \item \textbf{Nennen} Sie vier zentrale Aufgaben der Lagerhaltung. \hfill (4~P.)
  \item \textbf{Erläutern} Sie die \emph{Sicherungsfunktion} und \emph{Bereitstellungsfunktion} des Lagers. \hfill (4~P.)
\end{enumerate}

\subsection*{Aufgabe~3 --- Lagerformen und Lagerarten (10 Punkte)}

\begin{enumerate}[label=(\alph*)]
  \item \textbf{Vergleichen} Sie Eigenlager und Fremdlager (je 2~Vor-/Nachteile). \hfill (4~P.)
  \item \textbf{Nennen} Sie vier Lagerarten nach Bauart und beschreiben Sie kurz das Hochregallager. \hfill (4~P.)
  \item \textbf{Beurteilen} Sie: Welche Lagerart eignet sich für einen Bürodienstleister mit geringem Lagerbedarf --- Eigenlager oder Fremdlager? Begründen Sie. \hfill (2~P.)
\end{enumerate}

\subsection*{Aufgabe~4 --- Lagerkosten (8 Punkte)}

\noindent\textbf{\color{bfwAccent}YouTube-Vorbereitung:}\enspace \href{https://www.youtube.com/watch?v=qCRQkE5NIwA}{Lagerkennzahlen einfach erklärt --- Grundlagen Lagerkosten} (\url{https://www.youtube.com/watch?v=qCRQkE5NIwA})

\medskip
\begin{enumerate}[label=(\alph*)]
  \item \textbf{Nennen} Sie vier Bestandteile der Lagerkosten. \hfill (4~P.)
  \item \textbf{Erläutern} Sie den Begriff \emph{Kapitalbindungskosten} und warum sie zu den Lagerkosten zählen. \hfill (4~P.)
\end{enumerate}

\subsection*{Aufgabe~5 --- Lagerkennzahlen rechnen (16 Punkte)}

\noindent\textbf{\color{bfwAccent}YouTube-Vorbereitung:}\enspace \href{https://www.youtube.com/watch?v=0D0wVujGios}{Lagerkennzahlen IHK-Prüfung Sommer 2019} (\url{https://www.youtube.com/watch?v=0D0wVujGios})

\medskip
Aus der Lagerbuchhaltung der \emph{Schreibwaren-Vertrieb GmbH} liegen folgende Daten vor:

\begin{itemize}[itemsep=0pt]
  \item Anfangsbestand 01.01.: 4\,000~Stück (à 5,00~€)
  \item Endbestand 31.12.: 6\,000~Stück
  \item Jahresverbrauch: 50\,000~Stück
  \item Lagerkostensatz: 8\,\% p.\,a.
\end{itemize}

\begin{enumerate}[label=(\alph*)]
  \item \textbf{Berechnen} Sie den \textbf{durchschnittlichen Lagerbestand} in Stück und in €. \hfill (3~P.)
  \item \textbf{Berechnen} Sie die \textbf{Lagerumschlagshäufigkeit} (LUH). \hfill (3~P.)
  \item \textbf{Berechnen} Sie die \textbf{durchschnittliche Lagerdauer} in Tagen (Jahr = 360~Tage). \hfill (3~P.)
  \item \textbf{Berechnen} Sie die \textbf{Lagerzinsen} pro Jahr. \hfill (3~P.)
  \item \textbf{Beurteilen} Sie: Eine LUH von 10 --- wie würden Sie diese im Vergleich zur Branche (LUH \texttildelow12) bewerten? \hfill (4~P.)
\end{enumerate}

\subsection*{Aufgabe~6 --- IHK-Wissensfragen (6 Punkte)}

\begin{enumerate}[label=(\alph*)]
  \item \textbf{Erläutern} Sie kurz: Was bedeutet eine hohe Lagerumschlagshäufigkeit für ein Unternehmen? \hfill (2~P.)
  \item \textbf{Beurteilen} Sie: Warum ist die \emph{unverzügliche} Mängelrüge nach §\,377 HGB so wichtig? \hfill (2~P.)
  \item \textbf{Nennen} Sie zwei Gründe, warum hohe Lagerbestände wirtschaftlich nachteilig sein können. \hfill (2~P.)
\end{enumerate}

\vspace{1em}
\noindent\textbf{Punkte gesamt: 60}

\newpage

\section{Lösungen}\label{loesungen}

\subsection*{Lösung Aufgabe~1}
\begin{loesung}
\textbf{(a)} Die Wareneingangskontrolle erfolgt in vier Schritten:
\begin{enumerate}[itemsep=0pt]
  \item \emph{Annahme}: Lieferung entgegennehmen, Lieferschein unterschreiben (mit Vorbehalt bei Verdacht).
  \item \emph{Äußere Kontrolle}: Anzahl der Pakete zählen, sichtbare Beschädigungen vor dem Fahrer protokollieren.
  \item \emph{Innere Kontrolle}: Auspacken und Inhalt prüfen (Menge, Beschaffenheit, Identität), Vergleich mit Lieferschein und Bestellung.
  \item \emph{Erfassung}: Buchung im Warenwirtschafts­system, Einlagerung am bestimmten Lagerplatz.
\end{enumerate}

\textbf{(b)} \emph{Offene Mängel} sind bei der üblichen Wareneingangsprüfung sofort
erkennbar (z.\,B.\ falsche Stückzahl, beschädigte Verpackung). \emph{Verdeckte Mängel}
zeigen sich erst später, oft im Gebrauch (z.\,B.\ defekte Mechanik, schlechter Druck).
Beide müssen unverzüglich gerügt werden --- offene unverzüglich nach Lieferung,
verdeckte unverzüglich nach Entdeckung.

\textbf{(c)} §\,377 HGB regelt die \emph{Rügeobliegenheit} beim Handelskauf (Kauf
zwischen Kaufleuten). Der Käufer ist verpflichtet, die Ware unverzüglich nach Ablieferung
zu untersuchen und festgestellte Mängel unverzüglich anzuzeigen. \emph{Folge bei
Versäumnis}: Die Ware gilt als genehmigt --- alle Mängelrechte sind verloren.

\textbf{(d)} Mustermängelrüge:

\begin{quote}\small
Schreibwaren-Vertrieb GmbH, Berlin \\
20.~Mai 2026

\medskip
Bürofix AG, München

\medskip
\textbf{Mängelrüge zu Lieferung Bestell-Nr.\ 2026-0512-001 vom 20.~Mai 2026}

\medskip
Sehr geehrte Damen und Herren,

\medskip
am 20.~Mai 2026 erhielten wir Ihre Lieferung über 100 Aktenordner DIN~A4 (blau, breit).
Bei der Wareneingangsprüfung haben wir folgende Mängel festgestellt:

\begin{itemize}[itemsep=0pt]
  \item Fehlmenge: nur 95 statt 100 Stück geliefert
  \item 8 Aktenordner sind verschmutzt und teilweise eingerissen
\end{itemize}

\medskip
Wir bitten Sie um \emph{Nacherfüllung}: Lieferung der fehlenden 5~Stück sowie
Austausch der 8 beschädigten Aktenordner bis spätestens \textbf{27.~Mai 2026}.

\medskip
Mit freundlichen Grüßen \\
Ihr Name --- Sachbearbeitung Wareneingang
\end{quote}
\end{loesung}

\subsection*{Lösung Aufgabe~2}
\begin{loesung}
\textbf{(a)} Vier zentrale Aufgaben der Lagerhaltung:
\begin{itemize}[itemsep=0pt]
  \item Sicherungsfunktion (Schutz gegen Lieferengpässe und Bedarfsspitzen)
  \item Bereitstellungsfunktion (Material steht für Produktion/Verkauf bereit)
  \item Überbrückungsfunktion (Zeitpuffer zwischen Beschaffung und Verbrauch)
  \item Spekulationsfunktion (Einkauf zu günstigem Zeitpunkt vor Preiserhöhungen)
\end{itemize}

\textbf{(b)}
\begin{itemize}[itemsep=0pt]
  \item \emph{Sicherungsfunktion}: Das Lager wirkt als Versicherung gegen Lieferunterbrechungen, Streiks oder unerwartete Bedarfsspitzen. Auch wenn der Lieferant ausfällt, kann die Produktion vorübergehend weiterlaufen.
  \item \emph{Bereitstellungsfunktion}: Material ist jederzeit verfügbar, ohne dass auf den Lieferanten gewartet werden muss. Das beschleunigt Produktionsprozesse und verbessert die Lieferfähigkeit.
\end{itemize}
\end{loesung}

\subsection*{Lösung Aufgabe~3}
\begin{loesung}
\textbf{(a)} Vergleich:

\begin{center}
\begin{tabular}{p{4.5cm} p{4.5cm} p{4.5cm}}
\toprule
\textbf{Kriterium} & \textbf{Eigenlager} & \textbf{Fremdlager}\\
\midrule
Pro & Volle Kontrolle, eigenes Know-how & Keine Investition, variable Kosten\\
Pro & Direkte Anpassung an Bedarf & Hohe Mengenflexibilität\\
Contra & Hohe Investition und Fixkosten & Eingeschränkte Kontrolle\\
Contra & Geringe Flexibilität bei Größe & Lieferantenabhängigkeit\\
\bottomrule
\end{tabular}
\end{center}

\textbf{(b)} Vier Lagerarten nach Bauart: \emph{Bodenlager} (Waren am Boden, ohne Regal),
\emph{Regallager} (Standardregale, manuell oder mit Stapler), \emph{Hochregallager}
(bis 50~m Höhe, automatisiert), \emph{Blocklager} (Paletten gestapelt ohne Regale).

Das \emph{Hochregallager} repräsentiert die Königsklasse der raumoptimierten Lagerhaltung.
Mit Regalhöhen bis zu 50~m bietet es enorme Stellplatzkapazität auf minimaler
Grundfläche. Die Regale sind fest in die Hallenkonstruktion integriert und werden
fast ausschließlich automatisiert bewirtschaftet (Regalbediengeräte). Hohe Investition,
typisch für große Distributionszentren.

\textbf{(c)} Für einen Bürodienstleister mit geringem Lagerbedarf eignet sich ein
\textbf{Fremdlager}: Die hohen Investitionen für ein eigenes Lager (Gebäude, Einrichtung,
Personal) wären wirtschaftlich nicht darstellbar. Ein Fremdlager bietet variable Kosten
und passt sich dem schwankenden Bedarf an. Die etwas geringere Kontrolle ist bei
Bürobedarf akzeptabel.
\end{loesung}

\subsection*{Lösung Aufgabe~4}
\begin{loesung}
\textbf{(a)} Vier Bestandteile der Lagerkosten:
\begin{itemize}[itemsep=0pt]
  \item Personalkosten (Lagermitarbeiter, Disposition)
  \item Raumkosten (Miete, Energie, Reinigung)
  \item Anlagekosten (Abschreibung Regale, Stapler, IT)
  \item Kapitalbindungskosten (Opportunitätskosten für gebundenes Geld)
  \item (zusätzlich: Versicherungskosten, Lagerrisikokosten)
\end{itemize}

\textbf{(b)} \emph{Kapitalbindungskosten} sind Opportunitätskosten, die entstehen, weil
Geld in Lagerbeständen gebunden ist und nicht produktiv investiert werden kann. Beispiel:
Ein Unternehmen hat 100\,000~€ in Lagerwaren gebunden. Stattdessen hätte das Geld auf
einem Tagesgeldkonto 3\,\% Zinsen erbracht --- 3\,000~€ pro Jahr ,,verloren''. Diese
3\,000~€ sind die Kapitalbindungskosten. Sie zählen zu den Lagerkosten, weil sie direkt
durch die Lagerhaltung entstehen, auch wenn keine echte Auszahlung erfolgt.
\end{loesung}

\subsection*{Lösung Aufgabe~5}
\begin{loesung}
\textbf{(a)} Durchschnittlicher Lagerbestand:
$$\varnothing\,\text{Lagerbestand (Stück)} = \frac{\text{AB} + \text{EB}}{2} = \frac{4\,000 + 6\,000}{2} = \textbf{5\,000 Stück}$$
$$\varnothing\,\text{Lagerwert (€)} = 5\,000 \times 5{,}00\,€ = \textbf{25\,000 €}$$

\textbf{(b)} Lagerumschlagshäufigkeit (LUH):
$$\text{LUH} = \frac{\text{Jahresverbrauch}}{\varnothing\,\text{Lagerbestand}} = \frac{50\,000}{5\,000} = \textbf{10}$$

Das Lager wird also \emph{10-mal pro Jahr} komplett umgeschlagen.

\textbf{(c)} Durchschnittliche Lagerdauer:
$$\varnothing\,\text{Lagerdauer} = \frac{360 \text{ Tage}}{\text{LUH}} = \frac{360}{10} = \textbf{36 Tage}$$

Eine Ware liegt also durchschnittlich 36~Tage im Lager.

\textbf{(d)} Lagerzinsen pro Jahr:
$$\text{Lagerzinsen} = \varnothing\,\text{Lagerwert} \times \text{Lagerzinssatz} = 25\,000\,€ \times 0{,}08 = \textbf{2\,000 €}$$

Die jährliche Kapitalbindung verursacht 2\,000~€ Opportunitätskosten.

\textbf{(e)} Eine LUH von 10 ist im Vergleich zur Branche (LUH \texttildelow12) etwas
\emph{unter dem Branchenschnitt}. Das bedeutet:
\begin{itemize}[itemsep=0pt]
  \item Die Schreibwaren-Vertrieb GmbH hält im Schnitt mehr Lagerbestand als nötig --- Verbesserungspotenzial bei der Kapitalbindung.
  \item Konkrete Maßnahmen: Bestellmengen reduzieren (Andler-Formel prüfen), häufiger bestellen, Sicherheitsbestand kritisch hinterfragen.
  \item Eine LUH-Steigerung auf 12 würde den Ø~Lagerwert auf $50\,000 \div 12 \approx 4\,167$~Stück bzw.\ ca.\ 20\,833~€ senken --- 4\,167~€ weniger Kapitalbindung, ca.\ 333~€ weniger Lagerzinsen pro Jahr.
\end{itemize}
\end{loesung}

\subsection*{Lösung Aufgabe~6}
\begin{loesung}
\textbf{(a)} Eine \emph{hohe Lagerumschlagshäufigkeit} bedeutet, dass das Lager
schnell umgeschlagen wird --- Material liegt nur kurz, Kapitalbindung ist gering, das
Geld bleibt produktiv. Wirtschaftlich vorteilhaft. Allerdings auch: Hohes Risiko bei
Lieferengpässen, da kein großer Puffer.

\textbf{(b)} Die \emph{unverzügliche Mängelrüge} ist wichtig, weil bei Versäumnis die
Ware nach §\,377 Abs.\,2 HGB als genehmigt gilt --- der Käufer verliert alle
Mängelrechte (Nacherfüllung, Rücktritt, Minderung, Schadensersatz). Auch wirtschaftlich
schädigt eine späte Rüge: Beweissicherung wird schwieriger, Lieferanten-Beziehung
wird belastet.

\textbf{(c)} Zwei Nachteile hoher Lagerbestände:
\begin{itemize}[itemsep=0pt]
  \item \emph{Kapitalbindung}: Geld liegt im Lager statt produktiv investiert zu werden.
  \item \emph{Lager- und Risikokosten}: Höhere Miete, Versicherung, Verderb, Veralten, Schwund.
\end{itemize}
\end{loesung}

\end{document}
